\subsection{Ma trận nào mô tả quan sát, ma trận nào mô tả trung bình?}
\label{sec:covariance_roles}
Trước khi dựng miền tin cậy, cần biết ma trận đang đo độ phân tán của
\emph{đại lượng nào}. Các ma trận dưới đây có cùng kích thước $p\times p$
nhưng không có cùng vai trò.
\begin{center}
\small
\begin{tabular}{L{.19\linewidth}L{.25\linewidth}L{.44\linewidth}}
\toprule
Ma trận & Đại lượng được mô tả & Ý nghĩa và cách sử dụng\\
\midrule
$\Sigma$ & Một quan sát $\mathbf X_j$ & Hiệp phương sai tổng thể;
 mô tả đám mây quan sát, thường chưa biết.\\[3pt]
$S$ & Ước lượng $\Sigma$ & Hiệp phương sai mẫu không chệch;
 phản ánh độ phân tán và liên hệ giữa các cột dữ liệu.\\[3pt]
$\Sigma/n$ & Trung bình mẫu $\Xbar$ & Hiệp phương sai thực của trung bình
 mẫu khi các quan sát iid.\\[3pt]
$S/n$ & Ước lượng $\Cov(\Xbar)$ & Mô tả sai số của trung bình;
 nghịch đảo của nó tạo thước đo trong $T^2$.\\[3pt]
$\widehat\Sigma$ & MLE của $\Sigma$ & Bằng $(n-1)S/n$, phục vụ hàm hợp
 lý; không phải hiệp phương sai của trung bình mẫu.\\
\bottomrule
\end{tabular}
\end{center}
\nguon{Chương 4, §4.3, Result 4.11, tr.~171--172; §4.4, (4-23), tr.~174}

\paragraph{Đường chéo: biến thiên của từng biến riêng.}
$s_{ii}$ ước lượng phương sai của biến thứ $i$; $s_{ii}/n$ ước lượng
phương sai của trung bình thành phần. Do đó sai số chuẩn là
$\sqrt{s_{ii}/n}$. Không dùng $\sqrt{s_{ii}}$ làm sai số chuẩn của
trung bình: nó mô tả độ phân tán của quan sát.

\paragraph{Ngoài đường chéo: cách hai biến biến thiên cùng nhau.}
$s_{ik}$ mang thông tin ghép cặp giữa hai biến. Tương quan mẫu là
$r_{ik}=s_{ik}/\sqrt{s_{ii}s_{kk}}$. Nếu chia toàn bộ ma trận cho $n$,
tương quan không đổi vì hệ số $1/n$ triệt tiêu ở tử và mẫu.
Lấy nhiều quan sát giúp ước lượng trung bình chính xác hơn, nhưng không
đồng nghĩa hai thành phần trung bình mẫu trở nên không tương quan.

\paragraph{Ma trận nghịch đảo: độ lệch theo hướng nào đáng chú ý?}
Hotelling dùng $T^2=\mathbf r'(S/n)^{-1}\mathbf r$ với
$\mathbf r=\xbar-\vmu_0$. Phép nghịch đảo hiệu chỉnh đồng thời đơn vị,
phương sai và hiệp phương sai. Sau khi xoay theo các vectơ riêng,
\[
 T^2=n\sum_{i=1}^p\frac{(\mathbf e_i'\mathbf r)^2}{\ell_i},
 \qquad S\mathbf e_i=\ell_i\mathbf e_i.
\]
Độ lệch theo hướng có phương sai nhỏ $\ell_i$ bị đánh trọng số lớn hơn.
Đó là lý do một điểm gần tâm theo khoảng cách thường vẫn có thể nằm
ngoài elip. Nếu một trị riêng gần 0, nghịch đảo rất nhạy với sai số dữ
liệu; nếu trị riêng bằng 0, công thức cổ điển không dùng được.

\hieu{Hiệp phương sai không chỉ quyết định khoảng riêng của từng biến.
Nó còn cho biết các \emph{cặp giá trị trung bình} nào phù hợp với nhau.
Miền elip giữ thông tin về các hướng đó; một hình chữ nhật dựng từ hai
khoảng thành phần chỉ đặt giới hạn riêng cho từng tọa độ.}
